\section{Evidence and Witness Model}\label{sec:evidence}
Every record carries an evidence class (observation, derived, inferred, proof) and an optional status (missing, invalid, uncertain) in the specification. Runs and reproducibility checks are observations; the changed-region record and the witness difference are derived; the sufficiency record is derived and marked uncertain. Records are content-addressed and appended to a hash-chained JSONL log whose chain is verified on read. A trace function walks a verdict to its witness or sufficiency record, then to the runs, observations and program versions behind it. Run and witness identifiers are stable across invocations.

Implemented status: the run-set commitment is a flat hash of run ids, not a Merkle tree; supersession (voiding a witness) is not implemented; only a hit-line set per run is stored, not a hashed event stream; the dependency and runtime byte-identity check covers the source file only (Fig.~\ref{fig:evidence}). Reproducibility of the chain was checked by the tests and by whole-pair re-runs (identical verdicts, levels, reasons and witness ids; Table~\ref{tab:k1}), not by an independent re-derivation.
\dotfig{evidence_witness_model}{Evidence record classes and the witness and sufficiency chains}{fig:evidence}
